% Source PDF page 106; manual page 2-77.
\subsection{Aerodynamic Variables}\label{sec:aerodynamic-variables}

The ballistic coefficient is
\begin{equation}
  \beta_c
  =\frac{W}{C_DS_{\mathrm{ref}}}
  =\frac{mg}{C_DS_{\mathrm{ref}}}.
  \label{eq:ballistic-coefficient}
\end{equation}
It is a useful reentry-vehicle design parameter because it measures resistance to
aerodynamic deceleration. A vehicle with a low ballistic coefficient, for example
$500~\mathrm{lbf/ft^2}$, decelerates rapidly and generally experiences less
aerodynamic heating. A vehicle with a high ballistic coefficient, for example
$4000~\mathrm{lbf/ft^2}$, retains velocity longer and generally experiences greater
heating.

The lift-to-drag ratio is
\begin{equation}
  \frac{L}{D}
  =\frac{C_LqS_{\mathrm{ref}}}{C_DqS_{\mathrm{ref}}}
  =\frac{C_L}{C_D}.
  \label{eq:lift-to-drag-ratio}
\end{equation}
It measures aerodynamic efficiency. An unpowered vehicle obtains its greatest glide
performance near the angle of attack that maximizes $L/D$, denoted
$\alpha_{L/D}$.

\begin{figure}[htbp]
  \centering
  \resizebox{0.54\linewidth}{!}{\begin{tikzpicture}[>=Latex,line cap=round,line join=round,scale=0.9]
  \draw[->,thick] (0,0)--(5.3,0) node[below] {angle of attack};
  \draw[->,thick] (0,0)--(0,3.5) node[left] {$L/D$};
  \draw[very thick] plot[smooth] coordinates
    {(0,0) (1.1,1.25) (2.1,2.35) (3.2,3.05) (4.0,2.65) (4.7,2.1)};
  \coordinate (m) at (3.2,3.05);
  \fill (m) circle (2.4pt);
  \draw[densely dashed] (3.2,0)--(m);
  \draw[->] (4.65,3.4)--(3.32,3.08) node[pos=0,above] {maximum $L/D$};
\end{tikzpicture}
}
  \caption{Lift-to-drag ratio.}
  \label{fig:lift-to-drag-ratio}
\end{figure}

The function \texttt{maxld} finds $\alpha_{L/D}$ with the golden-section search
described in \cref{sec:search-methods}. The method is used because it converges
even when the tabulated $L/D$ function contains discontinuities.
